% called by main.tex
%
\chapter{Reproducibilidad}
\label{ch::reproducibilidad}

\section{Enlace a repositorio GitHub}

El código, los registros de resultados y los artefactos públicos de auditoría se mantienen en
un repositorio:

\begin{center}
\url{https://github.com/marcmaldonadolorca/polymarket-btc-microstructure}
\end{center}

\subsection{Estructura del repositorio}

El repositorio se organiza en los siguientes directorios:

\begin{description}[align=right,labelwidth=2.6cm]
\item [\texttt{config/}] Umbrales congelados de la política y configuraciones del arco.
\item [\texttt{docs/}] Documentos canónicos del proceso, renumerados de forma secuencial.
\item [\texttt{notebooks/}] Catorce cuadernos pedagógicos (00--13): el arco ejecutado
(EDA, particiones, objetivo, línea base, latencia, secuencias, corrección del \emph{score}
adverso, especialista y auditoría del ledger) más los cuadernos autocontenidos sobre CSV
públicos (EDA de series temporales, baselines clásicos, regímenes, economía maker y
fundacionales); cada uno declara qué pregunta responde, qué entradas usa, qué produce y su
conclusión, y enlaza su fila de \texttt{results/key\_results.csv}.
\item [\texttt{scripts/}] Puntos de entrada reproducibles de cada experimento.
\item [\texttt{src/}] Paquete con la lógica compartida (datos, características, modelos,
evaluación), con pruebas automatizadas.
\item [\texttt{results/}] Tablas de resultados y ledger anonimizado de 754 unidades de acción.
\item [\texttt{sql/}] Esquema de la consulta base de entrenamiento.
\item [\texttt{reports/}] Esta memoria (fuente LaTeX).
\end{description}

\subsection{Cómo reproducir}

El entorno se prepara con un entorno virtual de Python y la instalación de las dependencias
declaradas. Las pruebas automatizadas verifican la lógica de política y sus guardas. El ledger
\texttt{results/final\_candidate\_actions\_anonymized.csv} permite recalcular las tablas
agregadas, los desgloses por régimen y la incertidumbre agrupada sin acceder al corpus
completo. El script que lo genera queda versionado para mantener trazabilidad, aunque su
ejecución integral requiere el workspace privado de investigación.

\subsection{El recolector, y por qué no se publica}

Conviene decirlo con claridad porque es el artefacto más sustancial de ingeniería del trabajo
y el único que estuvo en producción: el sistema desplegado no fue nunca la política, fue el
\textbf{recolector}. Un servicio que corrió de forma continua durante meses sobre un
ordenador de placa reducida, alineando cuatro fuentes —libro y operaciones del contrato,
\emph{spot}, futuro perpetuo y oráculo— sobre una malla fija de dos segundos, con control de
calidad por sesión y rechazo de las que no lo superan. Todo lo que este trabajo puede afirmar
descansa en él, y las incidencias que acotaron sus ventanas —una reanudación tardía, una
desconexión del disco de almacenamiento— están reportadas en el Capítulo~\ref{ch::etica}
precisamente porque son propiedades del instrumento y no del análisis.

Su código \textbf{no se publica}, y el motivo no es el tamaño. Está acoplado a la
infraestructura privada donde vive —rutas, servicio, almacenamiento y la máquina concreta que
lo aloja— y publicarlo exigiría un trabajo de saneamiento que no se ha hecho: separar la
lógica de captura de la configuración del entorno y auditar que no arrastra nada del sistema
que lo hospeda. Publicar sin ese paso sería peor que no publicar.

Lo que sí se publica es lo que hace falta para reconstruirlo: el \textbf{esquema SQL} de la
base de captura, que define exactamente qué se guarda y con qué granularidad, y el
\emph{pipeline} de procesamiento posterior. Con eso, y con las dos APIs públicas de solo
lectura que el recolector consume —el catálogo de contratos y el libro de órdenes—, un
tercero puede escribir su propio recolector equivalente. No es lo mismo que entregar el
binario, y esa diferencia se declara aquí en vez de dejarla implícita: la reproducción
posible es la del \emph{método}, no la del \emph{corpus}.

\subsection{Datos no publicados}

El conjunto completo (del orden de 2,5 millones de filas \emph{cross-venue}, más la base de
captura) no se publica por tamaño. En su lugar se incluyen el esquema SQL, un ledger
anonimizado suficiente para reproducir las cifras finales y el código del \emph{pipeline}.
Esto permite una auditoría parcial, no la reproducción bit a bit del entrenamiento. Declarar
esa frontera de forma explícita es parte de la trazabilidad y reduce la deuda técnica propia de
estos sistemas~\cite{sculley2015debt}.
